% ═══════════════════════════════════════════════════════════════════════════════
\section{Der Intermediate Query Tree (IQT)}
\label{sec:iqt}
% ═══════════════════════════════════════════════════════════════════════════════

\subsection{Überblick und Architektur}
\label{subsec:iqt_uebersicht}

Das Herzstück des UQTL-Standards ist der \textbf{Intermediate Query Tree
(IQT)}: eine sprachunabhängige Baumstruktur, die jeden Query-Ausdruck aus
Python, Java und C\# einheitlich repräsentiert. Der IQT fungiert als
\emph{Pivot-Sprache} in der Übersetzungspipeline.

\begin{sidewaysfigure}
\centering
\begin{tikzpicture}[
  box/.style={draw, rounded corners=4pt, minimum width=2.8cm,
    minimum height=0.8cm, align=center, font=\small},
  source/.style={box, fill=blue!15},
  iqt/.style={box, fill=orange!25, minimum width=3.5cm, minimum height=1cm,
    font=\small\bfseries},
  target/.style={box, fill=green!20},
  arr/.style={-Stealth, thick},
  dbl/.style={Stealth-Stealth, thick, dashed},
]
  % Quellsprachen
  \node[source] (py)   at (0,4)   {Python\\SQLAlchemy};
  \node[source] (java) at (0,2)   {Java\\Streams / JPA};
  \node[source] (cs)   at (0,0)   {C\# / LINQ};

  % Parser
  \node[box, fill=yellow!20] (pp) at (3.5,4) {Python\\Parser};
  \node[box, fill=yellow!20] (jp) at (3.5,2) {Java\\Parser};
  \node[box, fill=yellow!20] (cp) at (3.5,0) {C\#\\Parser};

  % IQT
  \node[iqt] (iqt) at (7,2) {IQT\\Intermediate\\Query Tree};

  % Normalisierung
  \node[box, fill=purple!15] (norm) at (10.5,2) {IQT\\Normalisierer};

  % SQL-Generatoren
  \node[target] (ansi)  at (14,4)   {ANSI SQL:2016};
  \node[target] (pg)    at (14,2.7) {PostgreSQL};
  \node[target] (mysql) at (14,1.4) {MySQL 8};
  \node[target] (ms)    at (14,0.1) {MSSQL / Oracle};

  % Pfeile
  \draw[arr] (py)   -- (pp);
  \draw[arr] (java) -- (jp);
  \draw[arr] (cs)   -- (cp);
  \draw[arr] (pp)   -- (iqt);
  \draw[arr] (jp)   -- (iqt);
  \draw[arr] (cp)   -- (iqt);
  \draw[arr] (iqt)  -- (norm);
  \draw[arr] (norm) -- (ansi);
  \draw[arr] (norm) -- (pg);
  \draw[arr] (norm) -- (mysql);
  \draw[arr] (norm) -- (ms);

  % Beschriftung
  \node[font=\tiny, text=gray] at (1.7,4.3)  {AST-Extraktion};
  \node[font=\tiny, text=gray] at (1.7,2.3)  {AST-Extraktion};
  \node[font=\tiny, text=gray] at (1.7,0.3)  {AST-Extraktion};
  \node[font=\tiny, text=gray] at (5.2,3)    {IQT-Aufbau};
  \node[font=\tiny, text=gray] at (8.7,2.3)  {Normalisierung};
  \node[font=\tiny, text=gray] at (12.2,3)   {Codegen};
\end{tikzpicture}
\caption{UQTL-Übersetzungspipeline: von Quellsprachen über IQT zum SQL-Zieldialekt}
\label{fig:pipeline}
\end{sidewaysfigure}

\subsection{IQT-Knotentypen}
\label{subsec:iqt_knoten}

Der IQT ist ein gerichteter azyklischer Baum (DAG) aus typisierten Knoten.
Jeder Knoten entspricht einem relationalen oder logischen Operator.

\begin{definition}[IQT-Knotentypen]
\label{def:iqt_knoten}
Die Menge der zulässigen IQT-Knotentypen $\mathcal{N}$ ist definiert als:
\begin{align*}
\mathcal{N} = \{
  &\mathtt{SCAN},\; \mathtt{FILTER},\; \mathtt{PROJECT},\; \\
  &\mathtt{JOIN},\; \mathtt{GROUP\_AGG},\; \mathtt{SORT},\; \\
  &\mathtt{LIMIT},\; \mathtt{WINDOW},\; \mathtt{SUBQUERY},\; \\
  &\mathtt{SET\_OP},\; \mathtt{WITH}\;\}
\end{align*}
\end{definition}

\subsubsection{Knoten-Spezifikation im Detail}

\textbf{SCAN-Knoten}
Der \texttt{SCAN}-Knoten repräsentiert den Basistabellenzugriff:

\begin{lstlisting}[style=sqlstyle, caption={IQT-Pseudonotation: SCAN},
  label=lst:iqt_scan]
SCAN {
  table    : QualifiedName   -- Schema + Tabellenname
  alias    : Identifier?     -- optionaler Alias
  hint     : IndexHint?      -- optionaler Index-Hint (Erweiterung)
}
\end{lstlisting}

\textbf{FILTER-Knoten}
\begin{lstlisting}[style=sqlstyle, caption={IQT-Pseudonotation: FILTER},
  label=lst:iqt_filter]
FILTER {
  source   : IQTNode         -- Eingaberelation
  predicate: Expr            -- boolescher Ausdruck
  position : {WHERE | HAVING}
}

Expr ::=
  | CompareExpr { left: Expr, op: {EQ|NEQ|LT|GT|LTE|GTE}, right: Expr }
  | LogicExpr   { op: {AND|OR|NOT}, operands: Expr+ }
  | InExpr      { value: Expr, list: Expr+ | SubQuery }
  | BetweenExpr { value: Expr, lo: Expr, hi: Expr }
  | LikeExpr    { value: Expr, pattern: StringLiteral, escape: Char? }
  | NullExpr    { value: Expr, negated: Bool }
  | ExistsExpr  { subquery: SubQuery, negated: Bool }
  | Literal     { type: DataType, value: Any }
  | ColumnRef   { table: Identifier?, column: Identifier }
  | FuncCall    { name: QualifiedName, args: Expr*, distinct: Bool }
\end{lstlisting}

\textbf{PROJECT-Knoten}
\begin{lstlisting}[style=sqlstyle, caption={IQT-Pseudonotation: PROJECT},
  label=lst:iqt_project]
PROJECT {
  source    : IQTNode
  columns   : ProjectionItem+
  distinct  : Bool
}

ProjectionItem ::=
  | Star       {}
  | Expression { expr: Expr, alias: Identifier? }
\end{lstlisting}

\textbf{JOIN-Knoten}
\begin{lstlisting}[style=sqlstyle, caption={IQT-Pseudonotation: JOIN},
  label=lst:iqt_join]
JOIN {
  left      : IQTNode
  right     : IQTNode
  type      : {INNER | LEFT_OUTER | RIGHT_OUTER | FULL_OUTER | CROSS}
  condition : JoinCondition
}

JoinCondition ::=
  | OnCondition    { predicate: Expr }
  | UsingCondition { columns: Identifier+ }
  | Natural        {}
\end{lstlisting}

\textbf{GROUP\_AGG-Knoten}
\begin{lstlisting}[style=sqlstyle, caption={IQT-Pseudonotation: GROUP\_AGG},
  label=lst:iqt_group]
GROUP_AGG {
  source     : IQTNode
  groupBy    : Expr*           -- leer = Gesamt-Aggregation
  aggregates : AggItem+
  having     : Expr?
}

AggItem ::= {
  func  : {COUNT|SUM|AVG|MIN|MAX|STDDEV|VAR|STRING_AGG|...}
  arg   : Expr          -- COUNT(*) => StarExpr
  distinct: Bool
  alias : Identifier?
}
\end{lstlisting}

\textbf{SORT-Knoten}
\begin{lstlisting}[style=sqlstyle, caption={IQT-Pseudonotation: SORT},
  label=lst:iqt_sort]
SORT {
  source  : IQTNode
  items   : SortItem+
}

SortItem ::= {
  expr      : Expr
  direction : {ASC | DESC}
  nulls     : {NULLS_FIRST | NULLS_LAST | DEFAULT}
}
\end{lstlisting}

\textbf{LIMIT-Knoten}
\begin{lstlisting}[style=sqlstyle, caption={IQT-Pseudonotation: LIMIT},
  label=lst:iqt_limit]
LIMIT {
  source : IQTNode
  count  : Expr | Null    -- NULL = kein Limit
  offset : Expr | Null    -- NULL = kein Offset
}
\end{lstlisting}

\textbf{WINDOW-Knoten}
\begin{lstlisting}[style=sqlstyle, caption={IQT-Pseudonotation: WINDOW},
  label=lst:iqt_window]
WINDOW {
  source       : IQTNode
  definitions  : WindowDef+
}

WindowDef ::= {
  alias        : Identifier
  partitionBy  : Expr*
  orderBy      : SortItem*
  frame        : WindowFrame?
  func         : WindowFunc
  resultAlias  : Identifier
}

WindowFrame ::= {
  unit  : {ROWS | RANGE}
  start : FrameBound
  end   : FrameBound
}

FrameBound ::=
  | UNBOUNDED_PRECEDING | CURRENT_ROW | UNBOUNDED_FOLLOWING
  | Preceding { n: Expr } | Following { n: Expr }
\end{lstlisting}

\subsection{Normalisierungsregeln des IQT}
\label{subsec:iqt_norm}

Um die eindeutige SQL-Ausgabe zu garantieren, werden IQT-Bäume vor der
SQL-Generierung normalisiert. Die Normalisierung umfasst folgende Schritte:

\begin{regel}[Prädikat-Normalisierung]
\label{regel:pred_norm}
Zusammengesetzte \texttt{AND}-Prädikate werden in einer kanonischen
Reihenfolge sortiert (lexikografisch nach Spaltenname, dann Operator):
\[
  p_1 \wedge p_2 \wedge \ldots \wedge p_n \;\rightarrow\;
  p_{\sigma(1)} \wedge p_{\sigma(2)} \wedge \ldots \wedge p_{\sigma(n)}
\]
mit $\sigma$ als lexikografischer Permutation.
\end{regel}

\begin{regel}[Alias-Normalisierung]
\label{regel:alias_norm}
Generierte Tabellenaliase folgen dem Schema \texttt{t1, t2, \ldots} in
der Reihenfolge des Auftretens im Baum (Pre-Order-Traversal).
\end{regel}

\begin{regel}[Boolean-Vereinfachung]
\label{regel:bool_simp}
Redundante boolesche Ausdrücke werden eliminiert:
\begin{align*}
  p \wedge \mathtt{TRUE} &\;\rightarrow\; p \\
  p \vee \mathtt{FALSE}  &\;\rightarrow\; p \\
  \neg\neg p             &\;\rightarrow\; p \\
  p \wedge p             &\;\rightarrow\; p
\end{align*}
\end{regel}

\begin{regel}[SELECT-Reihenfolge]
\label{regel:select_order}
Explizit genannte Projekti\-ons\-spalten werden in der Reihenfolge ihrer
Deklaration ausgegeben. \texttt{SELECT *} wird nur bei echten Stern-Projektionen
verwendet; implizite Stern-Projektionen werden beibehalten.
\end{regel}

\begin{regel}[JOIN-Kanonisierung]
\label{regel:join_canon}
Mehrfach-Joins werden left-assoziativ normiert.
\texttt{CROSS JOIN} ohne \texttt{ON}-Klausel wird nur verwendet, wenn
explizit kein Prädikat vorliegt; sonst wird \texttt{INNER JOIN ... ON}
bevorzugt.
\end{regel}
